% ID: SEL-D0034-OA-E02
% ESTADO: cotejada. Segunda pasada visual y borrador compilado revisados.
% Fuente: Selectividad/matematicas_aplicadas/titular_septiembre_1_MATEMATICAS_APLICADAS.pdf
% SHA-256: c637bbde5c26c0b35b59202bd832ddc97990d3845365ce483fbf2a49832f73fe
% Página: 1; opción A; ejercicio original 2.

De una función continua y derivable, $f$, se sabe que la gráfica de la función derivada, $f'$, es una parábola que pasa por los puntos $(-1,0)$ y $(3,0)$, y que tiene su vértice en el punto $(1,-2)$.
\begin{enumerate}[label=\alph*)]
\item \textbf{(1.5 puntos)} Determine los intervalos de crecimiento y decrecimiento de la función $f$, así como la existencia de extremos.
\item \textbf{(1 punto)} Si $f(1)=2$, encuentre la ecuación de la recta tangente a la gráfica de la función $f$ en el punto de abscisa $x=1$.
\end{enumerate}
